\documentclass[conference]{IEEEtran}
\IEEEoverridecommandlockouts
\usepackage{cite}
\usepackage{amsmath,amssymb,amsfonts}
\usepackage{algorithmic}
\usepackage{graphicx}
\usepackage{textcomp}
\usepackage{xcolor}
\usepackage{url}
\usepackage{booktabs}
\usepackage{hyperref}
\def\BibTeX{{\rm B\kern-.05em{\sc i\kern-.025em b}\kern-.08em
    T\kern-.1667em\lower.7ex\hbox{E}\kern-.125emX}}

\begin{document}

\title{Benchmarking Rule-Based, Classical Machine Learning,\\
and Deep Learning Anomaly Detectors\\
in ETL Data Quality Pipelines}

\author{\IEEEauthorblockN{Theo Bigand}
\IEEEauthorblockA{\textit{ESILV --- M1 Data \& Artificial Intelligence} \\
\textit{Léonard de Vinci Pôle Universitaire}\\
Courbevoie, France \\
theobigand@gmail.com}
}

\maketitle

\begin{abstract}
Modern data pipelines increasingly depend on automated data quality
controls, yet Extract--Transform--Load (ETL) practitioners still rely
mostly on hand-crafted declarative rules. Recent work has started to
introduce machine learning (ML) anomaly detectors inside ETL workflows,
but systematic, apples-to-apples comparisons between rule-based,
classical ML, and deep learning approaches remain scarce---especially
when evaluated within a realistic pipeline and broken down by the type
of anomaly being injected. This paper presents an open-source ETL
pipeline integrating three families of anomaly detectors: statistical
rules (Z-score, IQR, Great Expectations), classical ML detectors
(Isolation Forest, LOF, One-Class SVM), and deep learning detectors
(Autoencoder, LSTM-Autoencoder). We evaluate the three families on
three datasets---Kaggle Credit Card Fraud, UCI Water Quality, and a
controlled synthetic dataset generated with Faker in which six distinct
anomaly types are injected at controlled rates. Beyond aggregate
precision, recall, F1, ROC-AUC and PR-AUC, we introduce a
per-anomaly-type heatmap that maps each method to the anomaly
categories it detects best. Our goal is to give practitioners a
reproducible, pipeline-integrated benchmark that answers the practical
question: \emph{which family of detectors should I use, for which kind
of data quality problem, at what operational cost?}
\end{abstract}

\begin{IEEEkeywords}
data quality, ETL, anomaly detection, benchmarking, machine learning,
Great Expectations, Isolation Forest, autoencoder
\end{IEEEkeywords}

\section{Introduction}

Extract--Transform--Load (ETL) pipelines are the backbone of modern data
infrastructure: they move raw operational data into analytical stores
where it feeds dashboards, reports, and machine-learning
applications~\cite{khatri2010data}. The value delivered downstream is
bounded by the quality of the data that reaches the destination, which
is why data quality has long been a first-class concern for data
engineers. Poor data quality has been estimated to cost organizations
an average of \$12.9M per year~\cite{gopa2025aiaugmented}, making
automated, continuous quality control a critical capability rather than
an optional one.

The dominant approach in industry remains rule-based: practitioners
hand-write declarative expectations---``column \texttt{amount} must be
positive'', ``column \texttt{email} must match a regex''---in
frameworks such as Great Expectations, Deequ, or
Soda~\cite{pipelinetoolssurvey2024}. These rules are auditable and
cheap to run, but they scale poorly with schema complexity, detect only
violations the author anticipated, and miss subtle multi-column or
distributional anomalies. To address these limits, a recent line of
work has proposed to complement or replace such rules with machine
learning anomaly detectors embedded directly inside the ETL
flow~\cite{gopa2025aiaugmented,marupaka2024mlpredictive,sokkula2024implementing}.
The existing literature, however, suffers from three gaps.

First, most published pipelines are conceptual
architectures~\cite{gopa2025aiaugmented,marupaka2024mlpredictive} that
do not ship with reproducible code or quantitative comparisons. When
experiments are reported, they typically exercise a single algorithm on
a single dataset---for instance, Isolation Forest on Credit Card
Fraud~\cite{sokkula2024implementing}, which achieved a near-random
ROC-AUC of 0.51 and thus offers limited insight into when ML-based
quality control is actually worth adopting.

Second, the anomaly detection community has produced excellent
generic benchmarks---ADBench~\cite{han2022adbench} compares 30
algorithms on 57 tabular datasets; PyOD~\cite{zhao2019pyod} provides a
unified implementation of dozens of detectors---but these benchmarks
evaluate algorithms \emph{in vacuum}, outside of an ETL context. They
do not consider the declarative rule-based baselines that dominate
production systems, and they do not measure end-to-end pipeline
integration costs (latency, memory, operational overhead).

Third, benchmarks typically report a single aggregated score per
dataset, hiding the fact that different families of detectors excel at
different kinds of anomalies. A Z-score detector is essentially unable
to flag categorical format violations; Great Expectations trivially
catches null-rate spikes but ignores multivariate
outliers~\cite{goldstein2016comparative}; autoencoders learn joint
distributions but can be fooled by concept drift. A practitioner
choosing a detector today lacks a fine-grained, actionable map from
anomaly type to recommended method.

\paragraph{Contributions.} This paper addresses these three gaps. We
present:
\begin{enumerate}
    \item \textbf{A reproducible ETL pipeline} integrating rule-based,
    classical ML, and deep learning anomaly detectors in a single
    unified framework, released as open-source code.
    \item \textbf{A systematic head-to-head benchmark} of three
    families---statistical rules (Z-score, IQR, Great Expectations),
    classical ML (Isolation Forest, LOF, One-Class SVM), and deep
    learning (Autoencoder, LSTM-Autoencoder)---on three datasets
    (Credit Card Fraud, UCI Water Quality, synthetic Faker), reporting
    precision, recall, F1, ROC-AUC, PR-AUC, and per-row latency.
    \item \textbf{A per-anomaly-type analysis.} Using a synthetic
    dataset in which six distinct anomaly categories (missing values,
    duplicates, univariate outliers, multivariate outliers, temporal
    drift, invalid formats) are injected at controlled rates, we
    produce a method-by-type heatmap that guides the choice of
    detector in production.
\end{enumerate}

The rest of the paper is organized as follows. Section~II reviews
related work on data quality in ETL and on anomaly detection
benchmarks. Section~III describes our pipeline architecture, the
detectors we compare, and the datasets and metrics we use. Section~IV
reports experimental results. Section~V discusses when rules, classical
ML, and deep learning are each the right tool, and highlights practical
trade-offs in latency and operational cost. Section~VI concludes and
outlines future work.

\section{Related Work}
\label{sec:related}

\emph{To be written. Planned subsections:}
\begin{itemize}
    \item Data quality in ETL pipelines (Khatri \& Brown
    \cite{khatri2010data}; Data Quality
    Survey~\cite{dataqualitysurvey2024})
    \item Rule-based and declarative approaches (Great Expectations,
    Deequ, Soda~\cite{pipelinetoolssurvey2024})
    \item Classical ML for anomaly detection (Isolation Forest, LOF,
    One-Class SVM; Goldstein \& Uchida~\cite{goldstein2016comparative};
    PyOD~\cite{zhao2019pyod})
    \item Deep learning (Autoencoder, LSTM-Autoencoder)
    \item Existing benchmarks: ADBench~\cite{han2022adbench} and why
    it does not cover rule-based baselines, ETL integration, or
    per-anomaly-type analysis
    \item ML in ETL pipelines: Gopa~\cite{gopa2025aiaugmented},
    Marupaka \& Rangineni~\cite{marupaka2024mlpredictive},
    Sokkula \& Vuppala~\cite{sokkula2024implementing} as illustrative
    limitations
\end{itemize}

\section{Methodology}
\label{sec:methodology}

\emph{To be written. Planned content: pipeline architecture, detector
families, datasets, metrics, experimental protocol.}

\section{Experiments and Results}
\label{sec:experiments}

\emph{To be written after the benchmark runs complete.}

\section{Discussion}
\label{sec:discussion}

\emph{To be written.}

\section{Conclusion}
\label{sec:conclusion}

\emph{To be written.}

\bibliographystyle{IEEEtran}
\bibliography{references}

\end{document}
